\begin{question}
	% 题目文本部分（可插入到article环境中）
	如图所示，半径$R=2.45\,\text{m}$的竖直光滑$\dfrac{1}{4}$圆弧轨道$AB$，其底端右侧是一个凹槽，凹槽右端连接一个半径$r=0.40\,\text{m}$的光滑圆轨道，轨道固定在竖直平面内，$D$与圆心连线与竖直方向成$\theta=37^\circ$角。一质量为$m=1\,\text{kg}$的滑板$Q$放置在凹槽内水平面上，其上表面刚好与$B$点和$C$点水平等高。开始时滑板静置在紧靠凹槽左端处，此时滑板右端与凹槽右端的距离$d=0.60\,\text{m}$。一质量也为$m$的小物块$P$（可视为质点）从$A$点由静止滑下，当物块滑至滑板右端时滑板恰好到达凹槽右端，撞击后滑板立即停止运动。已知物块与滑板间的动摩擦因数$\mu=0.75$，其余接触面的摩擦均可忽略不计，取重力加速度$g=10\,\text{m/s}^2$。求：
	\begin{enumerate}
		\item 滑块滑至$B$点时对$B$点的压力；
		\item 滑板的长度$l$；
		\item 请通过计算说明滑块滑上$CD$轨道后能否从$D$点飞出？
	\end{enumerate}
	
	\begin{tikzpicture}[scale=2] % 缩放比例，可按需调整
		
		% ---------- 1. 定义关键坐标 ----------
		% l=1.6m
		\coordinate (O1) at (0,0);       % 大圆圆心O₁
		\coordinate (A) at (-2.45,0);       % 滑块初始位置A（O₁左侧，距离R=3）
		\coordinate (B) at (0,-2.45);       % 大圆弧终点B（O₁正下方，距离R=3）
		\coordinate (O2) at (2.2,-2.05);      % 小圆圆心O₂（B右侧，水平距离l+d=6）
		\coordinate (C) at (2.2,-2.45);       % 小圆最低点C（O₂正下方，距离r=1）
		\coordinate (D) at ({2.2+0.4*sin(37)}, {-2.45+0.4+0.4*cos(37)}); % 小圆上一点D（角度45°）
		\coordinate (BS) at (0,-3);
		\coordinate (BQ) at (1.6,-3); 		
		\coordinate (BE) at (2.2,-3); 
		% ---------- 2. 绘制几何图形 ----------
		% 四分之一圆弧（A→B）：以O₁为圆心，R=3，从180°到270°
		\draw[thick] (A) arc[start angle=180, end angle=270, radius=2.45];
		
		% 滑块P：在A点绘制小矩形
		\filldraw[black] (A) rectangle ++(0.15,-0.15);
		
		% 水平木板Q：从B右侧开始，长度l=1.6，高度0.2
		\filldraw[brown] (B) rectangle ++(1.6,-0.2);
		
		% 右侧半圆轨道（C→D）：以O₂为圆心，r=1，从90°到-90°（顺时针）
		\draw[thick] (C) arc[start angle=-90, end angle=53, radius=0.4];
		
		% ---------- 3. 绘制辅助虚线 ----------
		\draw[dashed] (A) -- (O1);   % A到O₁的水平虚线
		\draw[dashed] (O1) -- (B);   % O₁到B的垂直虚线
		\draw[dashed] (O2) -- (D);   % O₂到D的斜向虚线
		
		% ---------- 4. 添加文字/数学标注 ----------
		\node[left] at (A) {$A$};               % 滑块A的标签
		\node[above right] at ($(A)+(0.15,-0.65)$) {$P$}; % 滑块P的标签
		\node[above] at ($(A)!0.5!(O1)$) {$R=2.45$};          % 大圆半径R的标签
		\node[right] at (O1) {$O_1$};           % 大圆圆心O₁的标签
		\node[above] at ($(0.5,-3)!0.5!(4.5,-3)$) {$Q$}; % 木板Q的标签
		%\node[below] at (BS) {$l$}; % 木板长度l的标签
		\node[below] at ($(4.5,-3.2)!0.5!(6,-3.2)$) {$d$}; % 木板到C的距离d的标签
		\node[right] at (C) {$C$};              % 小圆最低点C的标签
		\node[right] at ($(O2)!0.5!(D)$) {$r=0.4$};          % 小圆半径r的标签
		\node[left] at (O2) {$O_2$};      % 小圆圆心O₂的标签
		\node[above right] at (D) {$D$};        % 小圆上点D的标签
		\node[below left]at(B){$B$};
		
		\draw[|<->|] (BS) -- node[below] {$l$} (BQ);
		\draw[|<->|] (BQ) -- node[below] {$d$} (BE);
		\draw[dashed](O2)--($(C)+(0,-0.2)$);
		\draw[thick]($(C)+(0,-0.2)$)--($(B)+(0,-0.2)$);
		\draw[dash dot] (O2)--++(0,1);
		\draw[thick] ($(O2)+(0.12,0.2)$) arc[start angle=53, end angle=90, radius=0.2]node[midway,above]{\footnotesize$\theta$};
	\end{tikzpicture}
	
	
\end{question}
\begin{solution}
	
	
\end{solution}